% (c)~2012 Dimitrios Vrettos -d.vrettos@gmail.com
% (c)~2014 Claudio.carboncinii -claudio.carboncini@gmail.com
% (c) 2015 Daniele Zambelli daniele.zambelli@gmail.com

\input{\folder sistemi2gr_grafici.tex}

\chapter{Sistemi di secondo grado}

\inicapitolo{%
\begin{itemize}[left=0mm, nosep]
\item sistemi di secondo grado;
\item sistemi simmetrici;
\item problemi.
\end{itemize}
}

\section{Sistemi di secondo grado numerici}
\label{sec:sist2gr_sisteminumerci}

Ricordiamo che un sistema di equazioni non è altro che l'insieme di più 
equazioni con le stesse incognite. L'insieme delle soluzioni è dato 
dall'intersezione degli insiemi delle soluzioni delle singole equazioni.

\begin{definizione}{}{}
\ind{grado di un sistema}
Il \emph{grado di un sistema di equazioni}, se le equazioni che formano il 
sistema sono costituite da polinomi, è dato dal prodotto dei gradi delle 
equazioni che lo compongono.
\end{definizione}

\begin{esempio}{}{}
Determinare il grado dei seguenti sistemi di equazioni

\begin{itemize}
\item \(\sistema{-2x+3y=4 \\3x+5y-2=0}\) 
entrambe le equazioni sono di primo grado; il sistema è di primo grado;
\item \(\sistema{2x-y=0 \\x^2+6y^2-9=0}\) la prima equazione è di primo grado,
la seconda di secondo grado; il sistema è di secondo grado;
\item 
\(\sistema{x^2+y^2=0 \\y=3x^2-2x+6=0}\) 
entrambe le equazioni sono di secondo grado; il sistema è di quarto grado.
\end{itemize}
\end{esempio}

I sistemi di secondo grado sono dunque composti da un'equazione di secondo 
grado e da una di primo grado.
Il modo più semplice per risolverli, è la sostituzione:
\ind{metodo di sostituzione}
dall'equazione di primo grado esplicitiamo il valore di un'incognita e lo 
sostituiamo a quell'incognita nell'altra equazione.
Incontriamo questi sistemi in particolare quando vogliamo trovare le
intersezioni tra una retta e una conica: la retta infatti viene 
rappresentata da un'equazione di primo grado e la conica da un'equazione 
di secondo grado.
Ne seguito vedremo alcuni esempi.

\begin{esempio}{}{}
Risolvere il seguente sistema: \quad 
\(\sistema{{2x-y=0}\\{x^2+6y^2-9=0}}\)

Utilizziamo il metodo di sostituzione che abbiamo già visto per i sistemi 
di primo grado.
\begin{itemize}[nosep, left=0mm]
\item Ricaviamo una delle due incognite dall'equazione di primo grado e 
sostituiamo nell'equazione di secondo grado:
\[\sistema{y=2x \\ x^2 +6 \cdot (2x)^2 -9 = 0} 
\Rightarrow\sistema{y=2x \\ x^2+24x^2-9=0} 
\Rightarrow \sistema{y=2x \\ 25x^2-9=0}\]
\item risolviamo l'equazione di secondo grado in una sola incognita. Questa 
equazione è detta equazione risolvente del sistema:
\(25x^2-9=0 \Rightarrow x_1 = -\dfrac 3 5 \sor x_2 = \dfrac 3 5\)
\item Si sostituiscono i valori trovati per la \( x \) nella equazione di 
primo grado per trovare i valori corrispondenti della \( y \). 
Le coppie 
\(\coppia{x_1}{y_1}\) e 
\(\coppia{x_2}{y_2}\) se ci sono, si dicono soluzioni del sistema.
\[\sistema{{y=2x} \\ {25x^2-9=0}} 
\Rightarrow \sistema{x_1=-\dfrac 3 5 \\
y_1=2\cdot \tonda{-\dfrac 3 5}=-\dfrac 6 5} \sor 
\sistema{x_2=+\dfrac 3 5 \\
y_2=2\cdot \tonda{\dfrac 3 5}=+\dfrac 6 5 }\] 
quindi con soluzioni: \quad 
\(\coppia{-\dfrac 3 5}{-\dfrac 6 5}  \sor \coppia{\dfrac 3 5}{\dfrac 6 5}\)
\end{itemize}

\affiancati{.52}{.46}{
Le soluzioni del sistema possono essere interpretate geometricamente come 
i punti di intersezione tra la retta rappresentata dall'equazione \(y=2x\) 
e la curva rappresentata dall'equazione \(x^2+6y^2=9\). 
Con qualsiasi software che disegni grafici inseriamo le due equazioni e 
otteniamo la seguente figura.
La curva rappresentata dalla seconda equazione è una ellisse; 
i punti \(A\) e \(B\), intersezione tra retta ed ellisse, corrispondono 
alle soluzioni del sistema.
}{
\immagine[.8]{Ellisse e una retta che la interseca.}{\ells}
}
\end{esempio}

\begin{esempio}{}{}
Calcola le intersezioni tra la retta di equazione \(x -y +2 = 0\) e la 
conica di equazione \(x^2 +y -3x -1 = 0\).

\medskip
I punti di intersezione sono rappresentati da coppie di valori che sono 
soluzioni sia della prima equazione sia della seconda.
Si ottengono quindi risolvendo il seguente sistema: 
\(\sistema{x -y +2 = 0 \\x^2 +y -3x -1 = 0}.\)

Isoliamo la \(y\) dell'equazione di primo grado e sostituiamo 
nell'equazione di secondo grado 
\[\sistema{y = x + 2 \\
x^2 +\tonda{x +2} -3x -1 = 0} 
\Rightarrow\sistema{y = x +2 \\
x^2 -2x +1 = 0}\]

\affiancati{.60}{.38}{
L' \emph{equazione risolvente del sistema} è:\\
\(x^2 -2x +1 = 0\) e ha il discriminante uguale a zero presenta quindi due 
soluzioni reali coincidenti: \(x_1 = x_2 = 1\) da cui si possono ricavare i 
corrispondenti valori di \(y\) ovviamente anche loro coincidenti: 
\[\sistema{x_{1,~2} = 1 \\ y_{1,~2} = 1 +2 = 3}\] 
Quindi possiamo dire che la retta e la parabola si intersecano in due punti 
coincidenti (o che hanno un'intersezione doppia): \quad
\(\punto{A}{1}{3}\).

Questo punto è detto punto di tangenza tra retta e parabola.
}{
\immagine[.8]{Parabola e retta tangente in A.}{\prbl}
}

\end{esempio}

\begin{esempio}{}{}
Trova le intersezioni tra la retta di equazione \(3x +4y = 12\) e la 
circonferenza di equazione \(x^2 +y^2 = 4\).
\medskip 
Le intersezioni avranno per coordinate le coppie di numeri soluzioni del 
sistema:  \(\sistema{3x +4y = 12 \\ x^2 +y^2 = 4}.\)

Isoliamo \(y\) nell'equazione di primo grado e sostituiamola nell'equazione 
di secondo grado 
\[\sistema{y=-\dfrac 3 4x+3 \\
x^2+\tonda{-\dfrac 3 4x+3}^2-4=0} 
\Rightarrow\sistema{y=-\dfrac 3 4x+3 \\
x^2+\dfrac 9{16}x^2-\dfrac 9 2x+9-4=0} 
% \Rightarrow\sistema{y=-\dfrac 3 4x+3 \\
% \dfrac{25}{16}x^2-\dfrac{9}{2}x+5=0}
\]

\affiancati{.51}{.47}{
Otteniamo l'equazione risolvente:\\ 
\(\dfrac{25}{16}x^2 -\dfrac{9}{2}x +5 = 0 \sRarrow
25x^2 -36x +80\) Studiamo il discriminante della formula ridotta: 
\(\dfrac{\Delta}{4} =256 -2000 = -1744\) osserviamo che è negativo, 
l'equazione non ha soluzioni reali e quindi anche il sistema di 
equazioni non ha soluzioni reali e si dice \emph{impossibile}.
}{
\immagine[.8]{Circonferenza e retta esterna.}{\crch}
}
\end{esempio}


\conclusione
Un sistema di secondo grado, con equazione risolvente di secondo grado, 
rappresenta sempre l'intersezione tra una retta e una curva di secondo 
grado (circonferenza, parabola, ellisse o iperbole). 
Le soluzioni del sistema rappresentano i punti di incontro tra retta e 
curva. 
In base al segno del discriminante dell'equazione risolvente abbiamo:
\begin{itemize}[nosep]
\item \(\Delta < 0\) il sistema non ha soluzioni reali, 
retta e curva non hanno punti in comune;
\item \(\Delta = 0\) le soluzioni del sistema sono le coordinate di due 
punti coincidenti;
\item \(\Delta > 0\) le soluzioni del sistema sono le coordinate di due 
punti distinti.
\end{itemize}
\immagine[.6]{Retta e ellisse secanti, retta e parabola tangenti, retta e 
circonferenza esterne.}
{\rettaconica}

\subsection{Alcuni casi particolari}
\label{subsec:sistemi2gr_casi_particolari}
\ind{sistema indeterminato}  \ind{sistema impossibile}
Ci sono alcune situazioni particolari: sistemi indeterminati, 
impossibili o con una sola soluzione.
Li affrontiamo partendo da alcuni esempi.

\begin{esempio}{}{}
Risolvere il sistema 
\(\sistema{x^2-y^2=0\\x+y=0}\).

\medskip
Isoliamo la \(y\) dell'equazione di primo grado e sostituiamo 
nell'equazione di secondo grado. 
\[\sistema{y=-x \\
x^2-(-x)^2=0}\ 
\Rightarrow\sistema{y=-x \\
x^2-x^2=0} 
\Rightarrow\sistema{y=-x\\
0=0}.\]

\affiancati{.51}{.47}{
L'\emph{equazione risolvente del sistema} in questo caso è una 
\emph{identità} 
(uguaglianza vera) e tutte le coppie formate da numeri opposti (la prima 
equazione ci vincola ad avere \(y=-x\) ) sono soluzioni del sistema: 
\(\forall k\in \mathbb{R}\Rightarrow {I.S.}=\coppia{k}{-k}\).
Il sistema ha infinite coppie di numeri reali che lo soddisfano e si dice 
\emph{indeterminato}.

% La curva di secondo grado è formata dalle due rette \(x+y=0\) e \(x-y=0\)
% e la seconda equazione rappresenta la retta che si sovrappone alla
% precedente. %\TODO vorrei chiarire ma non so se riesco
Il sistema di partenza è dato da una curva di secondo grado, che corrisponde al
prodotto delle due diagonali \(x+y=0\) e \(x-y=0\), e da un'equazione lineare,
che corrisponde alla diagonale discendente.
}{
\immagine[.8]{Le due diagonali del primo e del secondo quadrante e la retta 
y=-x.}
{\iprba}
}
\end{esempio}

\begin{esempio}{}{}
Risolvere il sistema \(\sistema{\dfrac 3 
2x+y=0\\x^2-y^2=4}.\)

\medskip
Isoliamo la \(y\) dell'equazione di primo grado e sostituiamo 
nell'equazione di secondo grado 
\[\sistema{y=-\dfrac 3 2x \\
x^2-\tonda{-\dfrac 3 2x}^2=4}
\Rightarrow \sistema{y=-\dfrac 3 2x \\ x^2-\dfrac 9 4x^2=4}
\Rightarrow \sistema{y=-\dfrac 3 2x \\ -\dfrac 5 4x^2=4}
\Rightarrow \sistema{y=-\dfrac 3 2x \\ x^2 = -\dfrac{16}{5}}\]

\affiancati{.51}{.47}{
L'equazione risolvente del sistema: \\
\(x^2 = -\dfrac{16}{5}\) \\
non ha soluzioni, quindi il sistema è \emph{impossibile}.

L'equazione di secondo grado rappresenta una curva formata da due tratti:
sono i due rami di un'iperbole. La seconda equazione rappresenta una retta.
Vediamo che curva e retta non hanno punti di intersezione.
}{
\immagine[.8]{Retta che non interseca un'iperbole.}{\iprbb}
}
\end{esempio}

\begin{esempio}{}{}
Risolvere il sistema 
\(\sistema{x^2-y^2=4\\-x+y=-1}\).

\medskip
Isoliamo la \(y\) dell'equazione di primo grado e sostituiamo 
nell'equazione di secondo grado 
\[\sistema{y=x-1 \\
x^2-(x-1)^2-4=0}
\Rightarrow \sistema{y=x-1 \\
x^2-x^2+2x-1-4=0}
\Rightarrow \sistema{y=x-1\\2x=5}\]

L'\emph{equazione risolvente del sistema} in questo caso è l'equazione di 
primo grado \(2x-5=0\), la cui soluzione è \(x=\dfrac 5 2\). Si sostituisce 
il valore trovato nell'altra equazione e troviamo la soluzione del sistema 
che in questo caso è unica: 
\affiancati{.51}{.47}{
\[\sistema{y=x-1 \\2x=5} 
\Rightarrow\sistema{x=\dfrac 5 2 \\[3mm]
y=\dfrac 5 2-1=\dfrac 3 2}\] 

L'equazione di secondo grado rappresenta una iperbole e 
l'altra equazione rappresenta una retta; 
vediamo che curva e retta hanno un solo punto di intersezione.
}{
\immagine[.8]{Retta tangente a un'iperbole.}{\iprbc}
}
\end{esempio}
%TODO non è meglio tracciare anche l'iperbole?
% \pagebreak %----------------------------------

\conclusione
Se l'equazione risolvente non è un'equazione di secondo grado:
\begin{itemize}[nosep]
\item se l'equazione risolvente è una uguaglianza vera, il sistema è 
indeterminato;
\item se l'equazione risolvente è una uguaglianza falsa il sistema è 
impossibile;
% \item se l'equazione risolvente è di primo grado determinata, da essa si 
% ricava il valore dell'incognita e si sostituisce tale valore nell'altra 
% equazione. 
% Il sistema ha una sola soluzione (in questo caso non si parla di due 
% soluzioni coincidenti, come nel caso precedente di \(\Delta =0\)).
\item se l'equazione risolvente è di primo grado il sistema ha una sola 
soluzione (in questo caso non si parla di due soluzioni coincidenti, 
come nel caso precedente di \(\Delta =0\)).
\end{itemize}
\ind{equazione risolvente di un sistema}

% \ovalbox{\risolvii \ref{ese:6.1}, \ref{ese:6.2}, \ref{ese:6.3}, 
% \ref{ese:6.4}, 
% \ref{ese:6.5}, \ref{ese:6.6}, \ref{ese:6.7}.}

% TODO: Da verificare, portare nelle disequazioni di secondo grado e poi 
% in matliberagraf.sty!!!

\newcommand{\parabolau}[3][x]{% Retta crescente
  \asseo{-3}{+3}{0}{#1}
  \tkzInit[xmin=-3,xmax=+3,ymin=-2,ymax=+2]
  \tkzFct[domain=-3.3:-1, ultra thick, blue!50!black]{.15*(x**2 -1)}
  \tkzFct[domain=-1:+1, ultra thick, red!50!black]{.15*(x**2 -1)}
  \tkzFct[domain=+1:+3.3, ultra thick, blue!50!black]{.15*(x**2 -1)}
  \draw [black] (-1, -.3) -- (-1, +.8) node [above] {$#2$};
  \draw [black] (+1, -.3) -- (+1, +.8) node [above] {$#3$};
}

% \NewDocumentCommand \tresegni{m s m s m}{
%   % Segno di un polinomio di primo grado.
%   % Se c'è l'asterisco lo zero viene reso con una crocetta.
%   % Sintassi: \segni[*]{primosegno}{secondosegno}
%   % Esempio di chiamata:
%   % \disegno{\tresegni*{-}{+}
%   % \disegno{\tresegni{-}{+}
%   % \disegno{\tresegni{-}*{+}
%   \node at (-2.5, .4) {$#1$};
%   \node at (0, .4) {$#3$};
%   \node at (+2.5, .4) {$#5$};
%   {\IfBooleanTF{#2}
%       {\crocetta{-1}{.4}}
%       {\cerchietto{-1}{.4}}}
%   {\IfBooleanTF{#4}
%       {\crocetta{+1}{.4}}
%       {\cerchietto{+1}{.4}}}
% }

\NewDocumentCommand{\segniparu}{O{x} s m s m}{% Retta crescente
  \disegno{
    \parabolau[#1]{#3}{#5}
  {\IfBooleanTF{#2}
      {{\IfBooleanTF{#4}
      {\tresegni{+}*{-}*{+}}
      {\tresegni{+}*{-}{+}}}}
      {{\IfBooleanTF{#4}
      {\tresegni{+}{-}*{+}}
      {\tresegni{+}{-}{+}}}}}
  }
}

\section{Sistemi di secondo grado letterali}
\label{sec:sist2gr_sistemiletterali}
\ind{sistema letterale di secondo grado}

Un sistema letterale si risolve come nel caso degli analoghi sistemi numerici.
Bisognerà, nell'equazione risolvente, discutere per quali valore del 
parametro si otterranno soluzioni reali. 
Vediamo un esempio.

\begin{esempio}{}{}
Discutere e risolvere il seguente sistema: 
\(\sistema{y -kx = -2\\y -x^2 = 2}\).

\medskip
Ricaviamo y dalla prima equazione e sostituiamo nella seconda equazione:
\[ \sistema{y = kx -2 \\ kx -2 -x^2 = 2}\ 
\Rightarrow\sistema{y = kx -2 \\ -x^2 +kx - 4 = 0}\ 
\Rightarrow\sistema{y = kx -2 \\ x^2 -kx +4 = 0}\]

Risolviamo l'equazione risolvente di secondo grado: \\[-.5em]
\[x^2 -kx +4 = 0 \sRarrow x_{1,~2} = \dfrac{+k \mp \sqrt{k^2-16}}{2}\]
Dobbiamo studiare il segno del discriminante 
calcolando gli zeri del polinomio con l'equazione \(k^2 -16 = 0\) e 
tracciando il grafico della parabola \(y = k^2 -16\):

\segnopara{k_{1} = -4;~k_{2} = +4}
           {y = k^2 -16}
           {\segniparu[k]{-4}{+4}} 

Quindi:
\begin{itemize}
\item se \(\Delta > 0\) cioè se \(-4 < k < +4\) il sistema non ha soluzioni 
reali;
\item se \(\Delta = 0\) il sistema ha due soluzioni reali e 
coincidenti:
\subitem per \(k = -4\): \quad \(s_{1,~2} = \coppia{-2}{+6}\) 
\subitem per \(k = +4\): \quad \(s_{1,~2} = \coppia{+2}{+6}\)
\item se \(\Delta > 0\) cioè se \(k<-4 \sor k>4\) il sistema avrà due 
soluzioni diverse: 

\medskip
\(
\sistema{x_1 = \dfrac{k-\sqrt{k^2-16}}{2} \\ [3mm]
         y_1 = kx_1 -2}
\sor 
\sistema{x_2 = \dfrac{k+\sqrt{k^2-16}}{2} \\ [3mm]
         y_2 = kx_2 -2}
\)
\end{itemize}
\end{esempio}

% \ovalbox{\risolvii \ref{ese:6.8}, \ref{ese:6.9}.}

\section{Sistemi frazionari}
\label{sec:sist2gr_sistemifrazionari}
\ind{sistema frazionario}

Si dice frazionario un sistema in cui almeno una delle equazioni che lo 
compongono presenta l'incognita al denominatore.
In questo caso, prima di affrontare la soluzione del sistema, bisogna 
studiare l'insieme di definizione delle frazioni.

\begin{esempio}{}{}
Risolvere il seguente sistema 
\(\sistema{2x-y=2 \\ \dfrac{x}{y+2}=\dfrac{x}{2y+5}}\).

\pagebreak %--------------------------------
\medskip
Determiniamo le condizioni di esistenza delle frazioni:\\ [-.5em]
\[\dfrac{x}{y+2}=\dfrac{x}{2y+5} \sRarrow 
\CE y \neq -2 \sand y \neq -\dfrac{5}{2}\]

Ora possiamo affrontare la soluzione del sistema. 
Trasformiamo l'equazione frazionaria nella corrispondente equazione intera:
\begin{align*}
\dfrac{x}{y+2}=\dfrac{x}{2y+5} \sRarrow & x \cdot (2y+5) = x \cdot (y+2)\\
 & 2{xy}+5x-{xy}-2x=0 \sRarrow xy+3x=0.
\end{align*}

Possiamo risolvere questo sistema con il metodo di sostituzione: 
\[\sistema{y = 2x -2 \\ xy +3x = 0}
\sRarrow \sistema{y = 2x -2 \\ x(2x -2) +3x = 0}
\sRarrow \sistema{y = 2x -2 \\ 2x^2 +x = 0}\]

\(2x^2 +x = 0\) è l'equazione risolvente e ha soluzioni: 
\(x_1 = 0 \sor x_2 = -\dfrac{1}{2}\).

\medskip
Calcoliamo i corrispondenti valori di \(y\): \quad 
\(\sistema{x_1=0 \\ y_1=-2} \sor 
\sistema{x_2=-\dfrac{1}{2} \\[3mm] y_2=-3}\)

\medskip
Confrontiamo ora i risultati ottenuti con le condizioni di esistenza:
\begin{itemize}
\item \(\coppia{0}{-2}\) è una Soluzione Non Accettabile;
\item \(\coppia{-\dfrac{1}{2}}{-3}\) è una Soluzione Accettabile.
\end{itemize}
\ind{soluzione accettabile o non accettabile}
\end{esempio}

% \ovalbox{\risolvii \ref{ese:6.10}, \ref{ese:6.11}}

\section{Sistemi di secondo grado in tre incognite}
\label{sec:sist2gr_sistemitrei}
\ind{sistema di 2° grado con più equazioni}

Quanto detto si può estendere ai sistemi di secondo grado di tre o più 
equazioni 
con altrettante incognite. Per risolvere uno di tali sistemi si cercherà, 
operando successive sostituzioni a partire dalle equazioni di primo 
grado, di ottenere un'equazione di secondo grado in una sola incognita 
(equazione risolvente del sistema).

A partire dalle eventuali soluzioni di tale equazione, si determineranno 
poi le soluzioni del sistema.

\begin{esempio}{}{}
Risolvere il sistema: \quad 
\(\sistema{2x +y -z = 0 \\ 3x +4y -2z = 1 \\ xy -y^2 +z -5y = 0}\).

\medskip
Isoliamo \(z\) dalla prima equazione, che è di primo grado, e 
sostituiamo il suo valore nelle altre due equazioni: 
{\small
\[\sistema{z=2x+y\\
3x+4y-2(2x+y)=1\\
{xy}-y^2+(2x+y)-5y=0} \Rightarrow
% \sistema{z=2x+y\\
% 3x+4y-4x-2y=1\\
% xy-y^2+2x-4y=0} \Rightarrow
\sistema{z=2x+y\\
-x+2y-1=0\\
xy-y^2+2x-4y=0}\]
} % end small

Ricaviamo \(x\) dalla seconda equazione e la sostituiamo nelle altre 
equazioni: 
\[ \sistema{x=2y-1\\z=2(2y-1)+y\\2y^2-y-y^2+4y-2-4y=0} 
\Rightarrow\sistema{x=2y-1\\z=5y-2\\y^2-y-2=0}\]

L'equazione \(y^2-y-2 = 0 \sRarrow \tonda{x-2} \tonda{x+1} = 0\) è 
l'equazione risolvente del sistema, le sue soluzioni sono 
\(y_1 = +2 \sor y_2 = -1\).

Sostituiamo i valori trovati per la \(y\) nelle altre equazioni per 
trovare i valori corrispondenti della \(x\) e della \(z\): 
\[ \sistema{x=2(2)-1 = 3 \\ y = 2 \\ z=5(2)-2 = 8} \sor 
\sistema{x=2(-1)-1 = -3 \\ y = -1 \\ z=5(-1)-2 = -7}\Rightarrow 
\terna{3}{2}{8} \sor \terna{-3}{-1}{-7}\]
\end{esempio}

% \ovalbox{\risolvii \ref{ese:6.12}, \ref{ese:6.13}}

\section{Sistemi simmetrici}
\label{sec:sistemi2gr_simmetrici}
\ind{sistema di 2° grado simmetrico}

Un sistema di due equazioni in due incognite si dice \emph{simmetrico} se 
non cambia scambiando le incognite.
Per esempio, nel sistema 
\[\sistema{x+y=1\\x^2+y^2+3xy+5=0}
\sstext{scambiando} x \sstext{con} y \srarrow 
\sistema{y+x = 1 \\ y^2+x^2+3yx+5 = 0}\] 
che è identico al precedente.

Le soluzioni del sistema sono:  
\(\sistema{x_1=-2\\y_1=3} \stext{ e } \sistema{x_2=3\\y_2=-2}\) \\
e come si può notare \(x\) e \(y\) vengono scambiate anche nella soluzione.

In generale se il sistema è simmetrico trovata una coppia soluzione 
\(\coppia{a}{b}\) 
l'altra è \(\coppia{b}{a}\).

\subsection{Sistema simmetrico fondamentale}
Il sistema simmetrico fondamentale è del tipo 
\(\sistema{{x+y=s}\\{xy=p}}\) esso risolve 
il problema di trovare due numeri, nota la loro somma e il loro prodotto.

Ricordiamo che nell'equazione di secondo grado \(x^2+bx+c=0\), la somma 
delle radici è \(-b\), mentre il prodotto è \(c\). 
Pertanto, basta risolvere la seguente equazione, \emph{detta equazione 
risolvente: } \(t^2-st+p=0\) con \(s=-b\) e \(p=c\).
\ind{sistema simmetrico fondamentale}

In base al segno del discriminante abbiamo:
\begin{itemize}[nosep]
\item \(\Delta >0\) l'equazione risolvente ha due soluzioni distinte 
\(t_1\) e \(t_2\), le soluzioni del sistema sono: 
\(\sistema{{x_1=t_1}\\{y_1=t_2}} \sor 
\sistema{{x_2=t_2}\\{y_2=t_1}}\)
\item \(\Delta =0\) l'equazione risolvente ha radici coincidenti 
\(t_1=t_2\), le soluzioni del sistema sono: 
\(\sistema{{x_1=t_1}\\{y_1=t_1}} \sor 
\sistema{{x_2=t_1}\\{y_2=t_1}}\)
\item \(\Delta <0\) l'equazione non ammette soluzioni reali. Il sistema è 
impossibile.
\end{itemize}

\begin{esempio}{}{}
Risolvere il seguente sistema 
\(\sistema{{x+y=5}\\{xy=4}}\).

\affiancati{.51}{.47}{
L'equazione risolvente è \(t^2-5t+4=0\) le cui soluzioni sono: \(t_1=1 \sor 
t_2=4\).

Le soluzioni del sistema sono le seguenti: \[ 
\sistema{{x_1=1}\\{y_1=4}} \sor 
\sistema{{x_2=4}\\{y_2=1}}. \]

Possiamo interpretare i risultati ottenuti nel piano cartesiano: la retta 
di equazione \(x+y=5\) interseca l'iperbole equilatera \({xy}=4\) nei due 
punti \(\punto{A}{1}{4}\) e \(\punto{B}{4}{1}\).
}{
\immagine[.8]{Retta che interseca un'iperbole in due punti distinti.}
{\iprbd}
}
\end{esempio}

\begin{esempio}{}{}
Risolvere il seguente sistema 
\(\sistema{{x+y=1}\\{xy=2}}\).

\affiancati{.51}{.47}{
L'equazione risolvente è \[ t^2-t+2=0 \] che ha il discriminante negativo 
e dunque non ha soluzioni reali. 
Il sistema è impossibile.

Possiamo interpretare i risultati ottenuti nel piano cartesiano: la retta 
di equazione \(x+y=1\) non interseca l'iperbole equilatera \({xy}=4\).
}{
\immagine[.8]{Retta esterna a un'iperbole.}{\iprbe}
}
\end{esempio}

\begin{esempio}{}{}
Risolvere il seguente sistema 
\(\sistema{{x+y=2}\\{xy=1}}\).

\affiancati{.51}{.47}{
L'equazione risolvente è \(t^2-2t+1=0\) le cui soluzioni sono: 
\(t_1=t_2=1\).

Il sistema ha due soluzioni coincidenti: 
\[\sistema{{x_1=1}\\{y_1=1}} \sor 
 \sistema{{x_2=1}\\{y_2=1}}. \]

Possiamo interpretare i risultati ottenuti nel piano cartesiano: la retta 
di equazione \(x+y=2\) è tangente all'iperbole equilatera \(xy=1\) nel 
punto \(\coppia{1}{1}\).
}{
\immagine[.8]{Retta tangente a un'iperbole.}{\iprbf}
}
\end{esempio}

% \ovalbox{\risolvii \ref{ese:6.14}, \ref{ese:6.15}, \ref{ese:6.16}, 
% \ref{ese:6.17}, \ref{ese:6.18}, \ref{ese:6.19}}
% 
\subsection{Sistemi simmetrici riconducibili al sistema simmetrico 
fondamentale}
In questa categoria rientrano i sistemi simmetrici che, mediante artifici, 
possono essere trasformati in sistemi simmetrici del tipo precedente.
\ind{sistema simmetrico fondamentale}

\begin{esempio}{}{}
Risolvere il sistema: \quad 
\(\sistema{{x+y=a}\\
{x^2+y^2+{bx}+{by}=c}}\)

\medskip
È possibile trasformare il sistema in un sistema simmetrico fondamentale.

Ricordando lo sviluppo del quadrato del binomio, 
\((x +y)^2 = x^2 +2xy +y^2\) possiamo ricavare l'identità: \quad 
\(x^2+y^2=(x+y)^2-2{xy}\). \\[1mm]
Il sistema può essere riscritto così: \quad
\(\sistema{{x+y=a}\\
{(x+y)^2-2{xy}+b(x+y)=c}}\)

Ma osservando la prima equazione, possiamo sostituire, nella seconda, 
tutte le occorrenze di \(x +y\) con \(a\):
\[\sistema{x+y = a \\ a^2-2{xy}+{ba} = c} \Rightarrow 
  \sistema{x+y = a \\ xy = \dfrac{a^2+{ab}-c} 2}\]
Che è un sistema equivalente a quello di partenza, ma scritto nella forma 
fondamentale.
\end{esempio}

\begin{esempio}{}{}
Risolvere il sistema: \quad 
\(\sistema{{x+y=7}\\{x^2+y^2=25}}\)

Utilizzando l'identità \(x^2+y^2 = (x+y)^2-2{xy}\), e poi sostituendo 
\(x +y\) con \(7\), il sistema può essere riscritto così: 
\[\sistema{x+y=7 \\ (x+y)^2-2{xy}=25} \Rightarrow 
  \sistema{x+y=7 \\ 49-2{xy}=25}
% \Rightarrow \sistema{{x+y=7}\\
% {-2{xy}=25-49}}\Rightarrow 
\sistema{x+y=7 \\ {xy}=12}\]

L'equazione risolvente è \(t^2-7t+12=0\) le cui soluzioni sono 
\(t_1=3 \sor t_2=4\).

Le soluzioni del sistema sono: \quad
\(\sistema{{x_1=3}\\
{y_1=4}} \sor \sistema{{x_2=4}\\
{y_2=3}}\)
\end{esempio}

\begin{esempio}{}{}
Risolvere il sistema \quad 
\(\sistema{-3x-3y=-5 \\ 2x^2+2y^2=10}\)

Dividendo per \(-3\) la prima equazione, per \(2\) la seconda e 
utilizzando l'identità \(x^2+y^2=(x+y)^2-2{xy}\) 
si ha: \quad 
\(\sistema{{x+y=\dfrac 5 3}\\
{x^2+y^2=5}}
\Rightarrow \sistema{{x+y=\dfrac 5 3}\\
{(x+y)^2-2{xy}=5}}
\Rightarrow \sistema{{x+y=\dfrac 5 3}\\
{\tonda{\dfrac 5 3}^2-2{xy}=5}}\)

Da cui si ottiene il sistema simmetrico: \quad 
\(\sistema{x+y = \dfrac{5}{3} \\ xy = -\dfrac{10}{9}}\)

L'equazione risolvente è \(t^2-\dfrac{5}{3}t-\dfrac{10}{9} = 0\) le cui 
soluzioni sono:\\ 
\(t_1 = \dfrac{5-\sqrt{65}}{6} \sor t_2 = \dfrac{5+\sqrt{65}}{6}\).

Le soluzioni del sistema sono: 
\(\sistema{{x_1=\dfrac{5-\sqrt{65}} 
6}\\{y_1=\dfrac{5+\sqrt{65}} 
6}} \sor \sistema{{x_2=\dfrac{5+\sqrt{65}} 
6}\\{y_2=\dfrac{5-\sqrt{65}} 6}}\)
\end{esempio}

% \ovalbox{\risolvii \ref{ese:6.20}, \ref{ese:6.21}, \ref{ese:6.22}, 
% \ref{ese:6.23}, \ref{ese:6.24}}

% \subsection{Sistemi non simmetrici riconducibili a sistemi simmetrici}
% Rientrano in questa classe i sistemi che, pur non essendo simmetrici, 
% possono 
% essere trasformati, mediante opportune sostituzioni, in sistemi 
% simmetrici. 
% Naturalmente questi sistemi si possono risolvere anche con la procedura 
% solita 
% di sostituzione per i sistemi di secondo grado.
% 
% \begin{esempio}{}{}
% Risolvere il sistema 
% \(\sistema{{x-y=8}\\{{xy}=-15}}\).
% 
% Mediante la sostituzione \(y'=-y\) otteniamo 
% \(\sistema{{x+y'=8}\\{xy'=15}}\) che è un 
% sistema 
% simmetrico fondamentale.
% 
% L'equazione risolvente è \(t^2-8t+15=0\) le cui soluzioni sono 
% \(t_1=3 \sor t_2=5\), 
% pertanto il sistema ha le soluzioni 
% \[\sistema{{x_1=3}\\
% {{y_1}'=5}} \sor 
% \sistema{{x_2=5}\\
% {{y_2}'=3}}.\] 
% Dall'uguaglianza \(y'=-y\Rightarrow y=-y'\) otteniamo le soluzioni del 
% sistema dato 
% \[\sistema{{x_1=3}\\
% {y_1=-5}} \sor 
% \sistema{{x_2=5}\\
% {y_2=-3}}.\]
% \end{esempio}
% 
% \begin{esempio}{}{}
% Risolvere il sistema 
% \(\sistema{{2x-3y=8}\\{{xy}=2}}\).
% 
% Mediante la sostituzione \(x'=2x\) e \(y'=-3y\) da cui \(x=\dfrac{x'} 2\) 
% e 
% \(y=-\dfrac{y'} 3\) otteniamo 
% \[\sistema{{x'+y'=8}\\
% {\dfrac{x'} 2\cdot \tonda{-\dfrac{y'} 3}=2}}
% \Rightarrow\sistema{{x'+y'=8}\\
% {x'y'=-12}}\] che è un sistema simmetrico fondamentale.
% 
% Risolviamo il sistema simmetrico 
% \(\sistema{{x'+y'=8}\\{x'y'=-12}}\) con la 
% procedura nota. L'equazione risolvente è \(t^2-8t-12=0\) le cui soluzioni 
% sono 
% \(t_{1,2}=4\mp 2\sqrt 3\) pertanto il sistema ha le soluzioni: 
% \[\sistema{{x_1'=4-2\sqrt 7}\\
% {y_1'=4+2\sqrt 7}} \sor 
% \sistema{{x_2'=4+2\sqrt 7}\\
% {y_2'=4-2\sqrt 7}}.\] 
% Dalle sostituzioni \(x=\dfrac{x'} 2\) e \(y=-\dfrac{y'} 3\) otteniamo le 
% soluzioni 
% del 
% sistema iniziale 
% \[\sistema{{x_1=\dfrac{4-2\sqrt 7} 2=2-\sqrt 7}\\
% {y_1=\dfrac{-4-2\sqrt 7} 3}} \sor 
% \sistema{{x_2=\dfrac{4+2\sqrt 7} 2=2+\sqrt 7}\\
% {y_2=\dfrac{-4+2\sqrt 7} 3}}.\]
% \paragraph{Procedura di sostituzione}
% Ricaviamo una delle due incognite dall'equazione di primo grado e 
% sostituiamola 
% nell'altra equazione \[ \sistema{{2x-3y=8}\\
% {{xy}=2}}\Rightarrow
% \sistema{{y=\dfrac{2x-8} 
% 3}\\{{xy}=2}}\Rightarrow 
% \sistema{{y=\dfrac{2x-8} 
% 3}\\{x\tonda{\dfrac{2x-8} 3}=2}}\Rightarrow 
% \sistema{{y=\dfrac{2x-8} 3}\\{2x^2-8x-6=0}}\]
% 
% Risolviamo l'equazione \(2x^2-8x-6=0\) avente come soluzioni 
% \(x_1=2-\sqrt 7 \sor x_2=2+\sqrt 7\). 
% Applicando la formula ridotta otteniamo: 
% \(x_1=2-\sqrt 7 \sor x_2=2+\sqrt 7\).
% Sostituiamo i valori trovati e ricaviamo i valori della \(y\): \[ 
% \sistema{{x_1=2-\sqrt 7}\\{y_1=\dfrac{-4-2\sqrt 7} 
% 3}} \sor \sistema{{x_2=2+\sqrt 
% 7}\\{y_2=\dfrac{-4+2\sqrt 7} 3}}. \]
% \end{esempio}
% 
% % \ovalbox{\risolvi \ref{ese:6.25}}

% \subsection{Sistemi simmetrici di grado superiore al secondo}
% Introduciamo le seguenti trasformazioni dette formule di Waring, dal 
% nome del matematico che le ha formulate per primo. 
% Con tali formule, si possono trasformare le potenze di un binomio in 
% relazioni tra somme e prodotti delle due variabili che lo compongono. 
% Indicate come \(s\) somma delle variabili e \
% (p\) il 
% loro prodotto queste sono le prime formule fino alla potenza quinta.
% \begin{itemize}[nosep]
% \item \( a^2+b^2=(a+b)^2-2{ab}=s^2-2p \)
% \item \( a^3+b^3=(a+b)^3-3a^2b-3ab^2=(a+b)^3-3{ab}(a+b)=s^3-3{ps} \)
% \item \( a^4+b^4=s^4-4{ps}^2+2p^2 \)
% \item \( a^5+b^5=s^5-5{ps}^3+5p^2s \).
% \end{itemize}
% 
% \begin{esempio}{}{}
% Risolvere il sistema 
% \(\sistema{{x+y=1}\\{x^3+y^3-2{xy}=3}}\).
% 
% Applicando l'identità \(x^3+y^3=(x+y)^3-3{xy}(x+y)\), il sistema può 
% essere 
% riscritto così: \[ 
% \sistema{{x+y=1}\\{(x+y)^3-3{xy}(x+y)-2{xy}=3}}\Rightarrow 
% \sistema{{x+y=1}\\{1-5{xy}=3}}\Rightarrow 
% \sistema{{x+y=1}\\{{xy}=-\dfrac 2 5}}.\]
% 
% Da cui l'equazione risolvente \(t^2-t-\dfrac 2 5=0\Rightarrow 
% 5t^2-5t-2=0\) 
% con 
% \(t_1=\dfrac{5-\sqrt{65}}{10}\) e \(t_2=\dfrac{5+\sqrt{65}}{10}\).
% Le soluzioni del sistema sono: 
% 
% 
% \(\tonda{\dfrac{5-\sqrt{65}}{10};
% \dfrac{5+\sqrt{65}}{10}},\tonda{\dfrac{5+\sqrt{
% 65}}{10}}{\dfrac{5-\sqrt{65}}{10}}\).
% \end{esempio}
% 
% \begin{esempio}{}{}
% Risolvere il sistema \(\sistema{{x+y=-1}\\{x^4+y^4=\dfrac 7 
% 2}}\).
% 
% Ricordando l'identità \(x^4+y^4=(x+y)^4-4{xy}(x+y)^2+2x^2y^2\), 
% il sistema può essere riscritto così: 
% \[\sistema{{x+y=-1}\\{(x+y)^4-4{xy}(x+y)^2+2x^2y^2=\dfrac 7 
% 2}} \Rightarrow 
% \sistema{{x+y=-1}\\{2x^2y^2-4{xy}-\dfrac 5 
% 2=0}}.\]
% 
% Introduciamo l'incognita ausiliaria \(u=xy\). 
% L'equazione \(2x^2y^2-4{xy}-\dfrac 5 2=0\) diventa 
% \(2u^2-4u-\dfrac 5 2=0\) che ha come soluzioni 
% \(u_1=-\dfrac 1 2 \sor u_2=\dfrac 5 2\Rightarrow {xy}=-\dfrac 1 2 \sor 
% {xy}=\dfrac 5 2\).
% 
% Il sistema assegnato è equivalente a due insiemi 
% \[\sistema{{x+y=-1}\\{{xy}=-\dfrac 1 2}}
%  \sor 
% \sistema{{x+y=-1}\\{{xy}=\dfrac 5 2}}\] 
% e dunque il suo insieme soluzione \(S\) si ottiene dall'unione 
% dell'insieme soluzione dei due sistemi~\(S=S_1\cup S_2\).
% 
% Il primo sistema 
% \[\sistema{{x+y=-1}\\{{xy}=-\dfrac 1 2}}\] 
% ha equazione risolvente \(t^2+t-\dfrac 1 2=0\) con 
% \[t_1=\dfrac{-1-\sqrt 3} 2\text{ e }t_2=\dfrac{-1+\sqrt 3} 2.\] 
% Il sistema ha soluzioni 
% \[\coppia{\dfrac{-1-\sqrt 3} 2}{\dfrac{-1+\sqrt 3} 
% 2},\coppia{\dfrac{-1+\sqrt 
% 3} 
% 2}{\dfrac{-1-\sqrt 3} 2}\] e quindi \[S_1=\left\{\coppia{\dfrac{-1-\sqrt 
% 3} 
% 2}{\dfrac{-1+\sqrt 3} 2},\coppia{\dfrac{-1+\sqrt 3} 2}{\dfrac{-1-\sqrt 3} 
% 2}\right\}\]
% 
% Il secondo sistema \(\sistema{{x+y=-1}\\{{xy}=\dfrac 5 
% 2}}\) ha equazione risolvente \(t^2+t+\dfrac 5 2=0\), 
% che 
% ha \(\Delta <0\), l'insieme soluzione è vuoto. Anche il sistema non ha 
% soluzioni 
% reali, quindi \(S_2=\emptyset \). L'insieme soluzione del sistema 
% assegnato 
% \(\sistema{{x+y=-1}\\
% {x^4+y^4=\dfrac 7 2}}\) è 
% \(S=S_1\cup \emptyset =S_1\).
% \end{esempio}
% 
% \ovalbox{\risolvii \ref{ese:6.26}, \ref{ese:6.27}, \ref{ese:6.28}, 
% \ref{ese:6.29}, \ref{ese:6.30}, \ref{ese:6.31}, \ref{ese:6.32}, 
% \ref{ese:6.33}, 
% \ref{ese:6.34}, \ref{ese:6.35}}
% 
% \section{Sistemi omogenei di quarto grado}
% Un sistema si dice omogeneo se le equazioni, con l'eccezione dei termini 
% noti, 
% hanno tutti i termini con lo stesso grado. I sistemi omogenei di quarto 
% grado 
% sono quindi nella forma:
% \begin{equation*}
% \sistema{{{ax}^2+{bxy}+{cy}^2=d}\\
% {a'x^2+b'{xy}+c'y^2=d'}}
% \end{equation*}
% \paragraph{Primo caso} \(d=0 \wedge d'=0\).
% 
% Il sistema si presenta nella forma 
% \(\sistema{{{ax}^2+{bxy}+{cy}^2=0}\\{a'x^2+b'{xy}+c'y^2=0}
% }\). 
% Un sistema di questo tipo ha sempre almeno la soluzione 
% nulla 
% \(\coppia{0}{0}\).
% 
% Per trovare le soluzioni del sistema poniamo \(y={tx}\) sostituendo 
% abbiamo: 
% \[ 
% 
% 
% \sistema{{{ax}^2+{btx}^2+{ct}^2x^2=0}\\
% {a'x^2+b'{tx}^2+c't^2x^2=0}
% } \text{ da cui } 
% 
% 
% \sistema{{x^2(a+{bt}+{ct}^2)=0}\\
% {x^2(a'+b't+c't^2)=0}}. \]
% 
% Supponendo \(x\neq 0\) e che quindi che può assumere tutti i valori 
% \(k\in \R_0\) 
% possiamo dividere le due equazioni per \(x^2\), otteniamo così due 
% equazioni 
% nell'incognita \(t\) che possiamo risolvere. 
% Se le due equazioni ammettono qualche 
% soluzione comune allora il sistema ammette infinite soluzioni. Le 
% soluzioni 
% sono 
% del tipo \(x=k\) e \(y={kt}\), dove \(t\) è la soluzione comune di cui 
% si è detto prima.
% 
% \begin{esempio}{}{}
% Risolvere il seguente sistema \(\sistema{x^2-3{xy}+2y^2=0 
% \\-x^2+5{xy}-6y^2=0 }\).
% 
% Applichiamo la sostituzione \(y={tx}\), il sistema diventa 
% \(\sistema{x^2-3{tx}^2+2t^2x^2=0 \\-x^2+5{tx}^2-6t^2x^2=0 
% }\).
% 
% Dividendo per \(x^2\) otteniamo \(\sistema{1-3t+2t^2=0 
% \\1-5t+6t^2=0 
% }\).
% 
% La prima equazione è risolta per \(t_1=1 \sor t_2=\dfrac 1 2\), mentre la 
% seconda 
% equazione è risolta per \(t_3=\dfrac 1 2 \sor t_4=\dfrac 1 3\). Le due 
% equazioni 
% hanno una radice in comune \(t=\dfrac 1 2\).
% 
% Pertanto oltre alla soluzione \(\coppia{0}{0}\) il sistema ammette infinite 
% soluzioni 
% che 
% possono essere scritte nella forma 
% \(\sistema{x=k \\y=\dfrac 1 2k }\).
% \end{esempio}
% 
% 
% \paragraph{Secondo caso}\(d=0 \wedge d'\neq 0\).
% 
% Il sistema si presenta nella forma 
% 
% \(\sistema{{{ax}^2+{bxy}+{cy}^2=0}\\
% {a'x^2+b'{xy}+c'y^2=d'}}\).
% 
% Ponendo \(y={tx}\) si ha 
% \(\sistema{{{ax}^2+{btx}^2+{ct}^2x^2=0}\\{a'x^2+b'{tx}
% ^2+c't^2x^2=d'}}\).
% 
% Dividendo per \(x^2\) la prima equazione si ha 
% 
% \(\sistema{{a+{bt}+{ct}^2=0}\\{x^2(a'+b't+c't^2)=d'}}\).
% 
% Si risolve la prima equazione nell'incognita \(t\) si sostituiscono i 
% valori 
% trovati nella seconda equazione e si ricavano i valori di \(x\) e di 
% seguito i valori di \(y\) con \(y={tx}\).
% % \newpage
% \begin{esempio}{}{}
% Risolvere il sistema \(\sistema{x^2-{xy}-6y^2=0 
% \\-x^2+2{xy}-3y^2=-6 }\).
% 
% Sostituendo \(y={tx}\) il sistema diventa 
% \[\sistema{1-t-6t^2=0 \\x^2(-1+2t-3t^2)=-6 
% }.\]
% 
% La prima equazione ha per soluzioni \(t_1=\dfrac 1 3\) e 
% \(t_2=-\dfrac 1 2\).
% 
% Sostituendo \(t=\dfrac 1 3\) nella seconda equazione si ha 
% \(x_{1,2}=\mp 3\) 
% e 
% sapendo che \(y={tx}\) si ottengono le coppie 
% 
% 
% \[\sistema{x_1=3\\y_1=
% 1} \sor\left\{\begin{array}{
% l}x_2=-3\\y_2=-1}.\]
% 
% Sostituendo \(t=-\dfrac 1 2\) si ha \(x_3=-\dfrac{2\sqrt 6}{11} \sor 
% x_4=\dfrac{2\sqrt 
% 6}{11}\) e sapendo che \(y={tx}\) si ottengono le coppie 
% \[\sistema{x_3=-\dfrac{2\sqrt 6}{11}\\y_3=\dfrac{\sqrt 
% 6}{11}} \sor\sistema{x_4=-2\dfrac{\sqrt 
% 6}{11}\\y_4=-\dfrac{\sqrt 6}{11}}.\]
% 
% L'insieme soluzione del sistema è 
% \[\IS=\left\{(x_1}{y_1),(x_2}{y_2),(x_3}{y_3),(x_4}{y_4)\right\}.\]
% \end{esempio}
% 
% 
% \paragraph{Terzo caso} \(d\neq 0 \wedge d'\neq 0\).
% 
% Il sistema si presenta nella forma 
% 
% \(\sistema{{{ax}^2+{bxy}+{cy}^2=d}\\
% {a'x^2+b'{xy}+c'y^2=d'} }\).
% 
% Ponendo \(y=tx\) si ha 
% \[\sistema{{x^2(a+{bt}+{ct}^2)=d}\\
% {x^2(a'+b't+c't^2)=d'}}\]
% Dividendo membro a membro le due equazioni, sotto la condizione \(x\neq 
% 0\wedge 
% a'+b't+c't^2\neq 0\) otteniamo 
% \begin{align*}
% &\dfrac{a+{bt}+{ct}^2}{a'+b't+c't^2}=\dfrac d{d'}\\
% \Rightarrow & d'(a+{bt}+{ct}^2)=d(a'+b't+c't^2)\\
% \Rightarrow & ({cd}'-c'd)t^2+({bd}'-b'd)t+{ad}'-a'd=0 
% \end{align*}
% che è una equazione di secondo grado nell'incognita \(t\).
% 
% Se l'equazione ha come soluzioni \(t_1\) e \(t_2\) dobbiamo poi 
% risolvere i sistemi 
% \[ \sistema{y=t_1x \\a'x^2+b'{xy}+c'y^2=d' 
% } \sor\sistema{y=t_2x \\
% a'x^2+b'{xy}+c'y^2=d' }\]
% % \newpage
% \begin{esempio}{}{}
% Risolvere il sistema \(\sistema{x^2+3{xy}-y^2=-68 
% \\-2x^2+{xy}+3y^2=88 }\).
% 
% Sostituendo \(y={tx}\) il sistema diventa 
% \(\sistema{x^2(1+3t-t^2)=-68 \\x^2(-2+t+3t^2)=88 
% }\).
% 
% Dividendo membro a membro con la condizione 
% \(x\neq 0\wedge 3t^2+t-2\neq 0\) 
% cioè 
% \(x\neq 0\), \(t\neq -1\) e \(t\neq \dfrac 2 3\) si ha 
% \(\dfrac{1+3t-t^2}{-2+t+3t^2}=-\dfrac{68}{88}\), da cui l'equazione 
% \(29t^2+83t-12=0\) 
% con soluzioni \(t_1=\dfrac 4{29} \sor t_2=-3\).
% 
% A questo punto dobbiamo risolvere i due sistemi:\[ 
% \sistema{y=\dfrac 4{29}x \\-2x^2+{xy}+3y^2=88 
% } \sor\sistema{y=-3x \\-2x^2+{xy}+3y^2=88 
% }. \]
% 
% Il primo sistema è impossibile, il secondo ha soluzioni 
% 
% 
% \[\sistema{x_1=-2\\y_1=6}
%  \sor\left\{\begin{array}
% {l}x_2=2\\y_2=-6}.\] 
% L'insieme soluzione del sistema è \(\IS=\{(-2}{6),(2}{-6)\}\).
% \end{esempio}
% 
% \ovalbox{\risolvii \ref{ese:6.36}, \ref{ese:6.37}, \ref{ese:6.38}, 
% \ref{ese:6.39}, \ref{ese:6.40}, \ref{ese:6.41}, \ref{ese:6.42}, 
% \ref{ese:6.43}, 
% \ref{ese:6.44}, \ref{ese:6.45}, \ref{ese:6.46}, \ref{ese:6.47}, 
% \ref{ese:6.48}.}
% 
\section{Problemi che si risolvono con sistemi di grado superiore al 
primo}
\ind{problemi e sistemi di grado superiore}
Considerare più variabili ci permette di facilitare il processo di 
traduzione in linguaggio matematico.

\begin{problema}{}{}
Il trapezio isoscele \({ABCD}\) è inscritto in una semicirconferenza di 
diametro \({AB}\) di misura \(25\munit{cm}\) determinare le misure dei 
lati del trapezio sapendo che il perimetro è \(62\munit{cm}\).
\end{problema}

\needspace{3\baselineskip}
\begin{soluzione}~

\vspace{-.5em}
\affiancati{.51}{.47}{
\emph{Dati}: 

% \(\begin{array}{l}\overline{AB}=25;~2p=62;~\\
% {AB}\parallel {DC};~{AD}\equiv {CB}\end{array}\)

\hspace{10mm}
\(\overline{AB}=25;\quad 2p=62;\) 

\hspace{10mm}
\({AB}\parallel {DC};\quad {AD}\equiv {CB}\)

\emph{Obiettivo}: 

% \(\begin{array}{l}\overline{CB}\) \(\overline{DC}\end{array}\)

\hspace{10mm}
\(\overline{CB};\quad \overline{DC}\)

\medskip
\emph{Dati impliciti}: 

\hspace{10mm}
\(\overline{{KO}}=\overline{CH};\quad \overline{CO}=\dfrac{25}{2};\) 

\hspace{10mm}
\(\overline{KC} = \dfrac{\overline{DC}}{2};\quad
\overline{HB} = \dfrac{25-y}{2};\) 

\hspace{10mm}
\(\widehat{CKO} = 90\grado;\quad \widehat{CHB} = 90\grado\)

}{
\immagine{Il trapezio isoscele iscritto in una semicirconferenza.}
{\trapd}
}

\emph{Incognite}: 

\hspace{10mm}
\(\overline{{CB}}=x\) \(\overline{{DC}}=y\)

\emph{Vincoli: } 

\hspace{10mm}
\(\sistema{0<x<\dfrac{25}{2}\sqrt{2}\\0<y<25}\)

\emph{Relazioni tra dati e incognite}:

\hspace{10mm}
\(\sistema{{y+2x+25=62}\\
{\tonda{\dfrac{25}{2}}^2-\tonda{\dfrac{y}{2}}^2=x^2-
\tonda{\dfrac{25-y}{2}}^2}} \Rightarrow 
\sistema{{y=-2x+37}\\{x^2-25x+150=0}}\)

\emph{Soluzioni: } 

\hspace{10mm}
\(\sistema{{x_1=15}\\{y_1=7}} \sor 
\sistema{{x_2=10}\\{y_2=17}}\)

\emph{Verifica: }
Entrambe le soluzioni sono accettabili.
\end{soluzione}

La soluzione si basa sulla equazione di primo grado 
\(y+2x+25=62\) 
che definisce il perimetro, sulla congruenza dei segmenti 
\(\overline{{KO}}\) 
e 
\(\overline{{CH}}\) 
facilmente in quanto entrambe sono distanza tra due 
rette parallele e l'applicazione del teorema di Pitagora ai triangoli 
\({CKB}\) e \({CHB}\) 
rettangoli per costruzione. Naturalmente tutte le informazioni ausiliare 
vanno dimostrate ma, data la loro facilità, le lasciamo al lettore.

Importante è impostare le condizioni sulle incognite che devono essere 
maggiori di \(0\) ma anche \(x<\dfrac{25}{2}\sqrt{2}\) perché il 
trapezio non diventi un triangolo e \(y<25\) perché la base minore sia 
realmente minore. 
L'ultimo passo consiste nella verifica delle soluzioni, che nel nostro 
caso sono entrambe accettabili. 
Si hanno dunque due trapezi inscritti in quella semicirconferenza 
che avranno il perimetro di \(62\munit{cm}\).

\begin{problema}{}{}
L'azienda Profit intende fare una ristrutturazione riducendo il numero 
degli operai. 
Oggi spende per gli operai (tutti con lo stesso stipendio) \(800€\) al 
giorno. 
Se si licenziassero \(5\) dipendenti e si riducesse lo stipendio di 
\(2€\) al giorno si avrebbe un risparmio giornaliero di \(200€\). 
Quanti sono gli operai attualmente occupati nell'azienda?
\end{problema}

\begin{soluzione}~

\vspace{-.5em}
\emph{Dati}:\( \begin{array}{l}
\text{spesa per salari al giorno}= 800\munit{\text{€}};\\
\text{riduzione salario giornaliero}= 2\munit{\text{€}};\\
\text{riduzione numero operai}= 5\text{ unità};\\
\text{risparmio dopo il licenziamento e la riduzione di stipendio}= 
200\munit{\text{€}}.
\end{array}\)

\emph{Obiettivo}: numero operai occupati prima della ristrutturazione

\emph{Incognite}:\( \begin{array}{l}
x = \text{numero operai prima della ristrutturazione};\\
y = \text{salario percepito da ogni operaio prima della 
ristrutturazione}.
\end{array}\)

\emph{Vincoli}: \(\sistema{x\in \N\\y\in \R^+}\)

\emph{Altre Informazioni}: \( \begin{array}{l}
\text{Numero operai dopo la ristrutturazione} = x-5;\\
\text{salario dopo la ristrutturazione} =y-2;\\
\text{spesa per stipendi dopo la ristrutturazione} 
=800-200=600\munit{\text{€}}.\end{array}\)

\emph{Relazioni tra dati e incognite}:\\[-1em]
\begin{equation*}
\sistema{{{xy}=800}\\{(x-5)(y-2)=600}}\Rightarrow 
\sistema{{{xy}=800}\\{{xy}-2x-5y+10=600}}\Rightarrow 
\sistema{{{xy}=800}\\{2x+5y=210}}
\end{equation*}

\emph{Soluzioni}: 
\(\sistema{{x_1=25}\\{y_1=32}} \sor 
\sistema{{x_2=80}\\{y_2=10}}\)

\emph{Verifica}: Entrambe le soluzioni sono accettabili.
\end{soluzione}

Naturalmente c'è una grande differenza tra percepire 
\(32\munit{\text{€/giorno}}\) 
di salario o \( 10\munit{\text{€/giorno}} \), come avere impiegati 
\( 25 \) o \( 80 \) operai. 
Il problema va meglio definito. Basterebbe per questo un vincolo che 
ci dice qual è la paga minima giornaliera di un operaio.

\begin{problema}{}{}
Un numero \(k\in \N\) è composto da tre cifre. Il prodotto delle tre 
cifre è \(42\). 
Se si scambia la cifra delle decine con quella delle centinaia si 
ottiene un numero che supera \(k\) di \(360\). 
Se si scambia la cifra della unità con quella delle centinaia si ottiene 
un numero minore di \(99\) rispetto al numero \(k\). 
Trovare \(k\).
\end{problema}

\begin{soluzione}~

\vspace{-.5em}
\emph{Dati}: \(\begin{array}{l}
\text{il numero } k \text{ è composto da tre cifre};\\
\text{prodotto delle tre cifre } = 42;\\
\text{scambiando la cifra delle decine con quella delle centinaia 
otteniamo}l=k+360; \\
\text{scambiando la cifra delle unità con quella delle centinaia 
otteniamo} m=k-99.\\
\end{array}\)

\emph{Obiettivo}: trovare il numero \(k\).

\emph{Incognite}: \(\begin{array}{l}
x =\text{ cifra che rappresenta il numero delle centinaia;}\\
y=\text{ cifra che rappresenta il numero delle decine;}\\
z=\text{ cifra che rappresenta il numero delle unità.}
\end{array}\)

\emph{Vincoli}: 
\(\sistema{x\in \{1,2,3,4,5,6,7,8,9\} \\
y,z\in \{0,1,2,3,4,5,6,7,8,9\}}\).

\emph{Altre Informazioni}: \(\begin{array}{l}
k=100x+10y+z;\\
l=100y+10x+z;\\
m=100z+10y+x.
\end{array}\)

\emph{Relazioni tra dati e incognite}: 
\[ \sistema{x\cdot y\cdot z=42 \\
100y+10x+z=100x+10y+z+360\\100z+10y+x=100x+10y+z-99}\Rightarrow 
\sistema{x \cdot y \cdot z = 42 \\ x-y = -4\\x-z = 1}\]

\emph{Soluzioni}: 
\(\sistema{x_1=3\\y_1=7\\z_1=2}\).

\emph{Verifica}: La soluzione soddisfa le condizioni il numero cercato è 
\(372\).
\end{soluzione}

% \vspazio
% \ovalbox{\risolvii \ref{ese:6.49}, \ref{ese:6.50}, \ref{ese:6.51}, 
%          \ref{ese:6.52}, \ref{ese:6.53}, \ref{ese:6.54}, \ref{ese:6.55}, 
%          \ref{ese:6.56}, \ref{ese:6.47}, \ref{ese:6.58}, \ref{ese:6.59}, 
%          \ref{ese:6.60}, \ref{ese:6.61}}
% 
% \vspazio
% \ovalbox{\ref{ese:6.62}, \ref{ese:6.63}, \ref{ese:6.64}, \ref{ese:6.65}, 
%          \ref{ese:6.66}, \ref{ese:6.67}, \ref{ese:6.68}, \ref{ese:6.69}, 
%          \ref{ese:6.70}, \ref{ese:6.71}}
% % >8
